\section{Introduction}%
\label{sec:introduction}

Model-Based Testing (MBT) offers a structured approach to verification, replacing traditional manual test design with the automated generation of test cases from formal behavioral specifications~\cite{taxonomy_mbt}. Traditional manual testing processes are frequently unstructured, difficult to reproduce, and heavily dependent on the subjective experience of individual engineers. By relying on explicit models that encode the intended behaviors of a System Under Test (SUT) or its operational environment, MBT provides a structured and verifiable foundation for quality assurance.

Within the broader spectrum of model-based verification, Model-Based Statistical Testing (MBST) also referred to as statistical usage testing is a framework that introduces a probabilistic framework designed to evaluate software quality from the perspective of its operational environment~\cite{whittaker1993markov,whittaker1994markov}. Originating from foundational work at the Software Quality Research Laboratory and pioneering research by Whittaker, Thomason, Poore, and Walton, MBST shifts the testing objective from absolute code coverage to statistical inference regarding expected field quality.

Instead of treating all execution paths with equal weight, MBST models the infinite population of potential system uses by constructing a formal \emph{usage model}. This model is structurally represented as a directed graph where the nodes signify discrete \emph{states-of-use} and the directed arcs represent transition stimuli triggered by users, hardware components, or environmental factors.

\begin{figure}[!t]
\centering
\begin{tikzpicture}[->, >=stealth', auto, semithick, node distance=3cm]
  % Define styles
  \tikzstyle{state}=[circle, draw, minimum size=1.2cm, inner sep=0pt, font=\small]
  \tikzstyle{terminal}=[rectangle, rounded corners, draw, minimum size=0.8cm, inner sep=4pt, font=\small]

  % Nodes
  \node[terminal]                (begin) {Begin};
  \node[state, right of=begin]   (A)     {State A};
  \node[state, right of=A]       (B)     {State B};
  \node[terminal, right of=B]    (end)   {End};

  % Edges
  \path (begin) edge node {} (A);
  \path (A)     edge [loop above] node {0.90} (A);
  \path (A)     edge [bend left]  node {0.10} (B);
  \path (B)     edge [bend left]  node {0.10} (A);
  \path (B)     edge [loop above] node {0.90} (B);
  \path (B)     edge node {} (end);
\end{tikzpicture}
\caption{A simple Markov chain usage model. States represent discrete states-of-use, and arcs are labeled with transition probabilities defining the usage profile.}%
\label{fig:markov_usage_model}
\end{figure}

Stochastic parameters are integrated into this transition graph by defining a \emph{usage profile}, which assigns specific transition probabilities to every permissible transition arc. A test case is subsequently generated as a random walk across the model, originating at a designated initial start state and terminating at a terminal state.

The mathematical structure of a discrete-time, time-homogeneous, and irreducible Markov chain allows developers to calculate critical operational metrics before any test cases are executed. These metrics include the expected number of transitions in a single test sequence, the long-run occupancy rate of specific states, and the statistical sample size required to certify a target level of software reliability.

Historically, the construction of these usage models has been a manual and labor-intensive process~\cite{llm_state_machine_frameworks}. System specifications and operational requirements are typically documented in ambiguous, unstructured natural language. Translating these narrative requirements into verified state-transition topologies and mathematically consistent probability matrices requires significant formal methods expertise. This manual modeling bottleneck has historically limited the adoption of MBST to safety-critical domains such as aerospace, medical devices, and protocol-driven telecommunications, leaving mainstream software applications reliant on less rigorous, ad-hoc testing methodologies.

The central question this paper addresses is therefore not merely one of model-construction automation, but of \emph{fault-detection consequence}: does automating usage-model construction with sufficient structural fidelity translate into meaningfully better fault-revealing test suites compared with traditional testing approaches?  Existing work on LLM-assisted state-machine generation~\cite{llm_state_machine_frameworks} has demonstrated syntactic extraction capability, but has not evaluated whether generated models preserve the statistical properties transition probability calibration, full transition coverage, operationally weighted path sampling that are necessary for MBST's fault-detection guarantees to hold.

\subsection*{Research Questions}
\label{subsec:rqs}

This paper is structured around four research questions:

\begin{enumerate}
    \item[\textbf{RQ1}] \textbf{Structural fidelity.}  Can automated neuro-symbolic model construction achieve the state and transition extraction accuracy ($F_1 \geq 0.90$) required for safety-critical test generation, without manual modelling effort?

    \item[\textbf{RQ2}] \textbf{Fault-detection coverage.}  Does a NeSy-MBST-generated usage model produce test suites with higher transition and state coverage than a purely neural baseline, and how does this translate into fault-detection path diversity?

    \item[\textbf{RQ3}] \textbf{Probabilistic calibration.}  Does the symbolic constraint optimization preserve the statistical properties of the usage model specifically, operationally weighted transition probabilities that determine how test effort is allocated across fault-revealing paths?

    \item[\textbf{RQ4}] \textbf{Component necessity.}  Which architectural components of NeSy-MBST (symbolic verification loop, convex optimizer, closed-loop feedback) are individually necessary for achieving the above, and which failure modes does each component address?
\end{enumerate}

\subsection*{Paper Structure}

Section~\ref{sec:literature_review} reviews the 20-paper landscape spanning MBT, statistical model checking, and probabilistic testing, with emphasis on empirical fault-detection evidence.  Section~\ref{sec:related_work} characterises the existing MBST tool landscape and LLM-based state-machine generation approaches.  Section~\ref{sec:bottlenecks} formalises the four structural bottlenecks that prevent purely neural approaches from guaranteeing MBST's fault-detection properties.  Section~\ref{sec:framework} presents the NeSy-MBST architecture and its five components.  Section~\ref{sec:mathematical_formulations} provides the mathematical foundations.  Section~\ref{sec:evaluation} presents the empirical evaluation addressing RQ1--RQ3.  Section~\ref{sec:eval:ablation} presents the ablation study addressing RQ4.  Section~\ref{sec:threats} discusses threats to validity.  Section~\ref{sec:conclusion} provides adoption recommendations and future directions.
